%% Thai CV - moderncv banking, Thai (Sarabun) with English technical terms
%% Compile with: xelatex -interaction=nonstopmode -halt-on-error <file>.tex
%% Font paths below are relative to cv/ (where /apply writes the output).

\documentclass[11pt,a4paper,sans]{moderncv}
\moderncvstyle{banking}
\moderncvcolor{blue}

\renewcommand*{\namefont}{\fontsize{34}{36}\bfseries\upshape}
\colorlet{firstnamecolor}{color1}
\colorlet{lastnamecolor}{color1}
\colorlet{namecolor}{color1}
\renewcommand*{\sectionstyle}[1]{{\sectionfont\color{color1}#1}}

% moderncv loads fontspec itself under xetex; these calls must come after
% \documentclass to win. The sans option routes the document through \sfdefault,
% so \setsansfont is the one that actually changes the body text.
\usepackage{fontspec}
\defaultfontfeatures{Ligatures=TeX}
\setmainfont[Path = ../cover_letters/OpenFonts/fonts/sarabun/, Extension = .ttf,
  UprightFont = *-Regular, BoldFont = *-Bold, ItalicFont = *-Italic,
  BoldItalicFont = *-BoldItalic]{Sarabun}
\setsansfont[Path = ../cover_letters/OpenFonts/fonts/sarabun/, Extension = .ttf,
  UprightFont = *-Regular, BoldFont = *-Bold, ItalicFont = *-Italic,
  BoldItalicFont = *-BoldItalic]{Sarabun}
\newfontfamily\thaifont[Path = ../cover_letters/OpenFonts/fonts/sarabun/, Extension = .ttf,
  UprightFont = *-Regular, BoldFont = *-Bold, ItalicFont = *-Italic,
  BoldItalicFont = *-BoldItalic]{Sarabun}
\newfontfamily\latinfont[Path = ../cover_letters/OpenFonts/fonts/lato/, Extension = .ttf,
  UprightFont = *-Reg, BoldFont = *-Bol, ItalicFont = *-RegIta,
  BoldItalicFont = *-BolIta]{Lato}

\usepackage{ucharclasses}
\setTransitionsForThai{\thaifont}{\latinfont}

% Thai is written without spaces between words; ICU line breaking plus a
% stretchable inter-character skip is what lets XeTeX break a Thai paragraph.
\XeTeXlinebreaklocale "th"
\XeTeXlinebreakskip = 0pt plus 1pt

\AtEndPreamble{\hypersetup{
    colorlinks=true,
    linkcolor=blue,
    filecolor=magenta,
    urlcolor=blue,
    pdftitle={[YOUR_NAME] - CV},
    pdfpagemode=UseNone,
}}
\usepackage[scale=0.77]{geometry}
\usepackage{needspace}

% Personal data
\name{[FIRST_NAME]}{[LAST_NAME]}
\address{[YOUR_ADDRESS]}{}{}
\phone[mobile]{[YOUR_PHONE]}
\email{[YOUR_EMAIL]}
\extrainfo{\href{[YOUR_LINKEDIN_URL]}{LinkedIn}, \href{[YOUR_GITHUB_URL]}{GitHub}}

% Local-employer variant only. A photo, date of birth and nationality are
% customary for Thai employers and unwelcome at MNCs - enable deliberately.
% \photo[64pt][0.4pt]{[YOUR_PHOTO_FILE]}
% \extrainfo{\href{[YOUR_LINKEDIN_URL]}{LinkedIn} \quad
%   วันเกิด: [YOUR_DOB] \quad สัญชาติ: [YOUR_NATIONALITY]}

\begin{document}

% fontawesome5 contact icons sit in the title block; a ucharclasses transition
% would re-select the text font over them, so keep transitions off until after.
\XeTeXinterchartokenstate=0
\makecvtitle
\XeTeXinterchartokenstate=1

% ============================================================
%     PROFILE STATEMENT
% ============================================================

\section{ประวัติส่วนตัว}
\vspace{1pt}
\small{[Write a 3-5 line profile statement in Thai, tailored to the role. Keep
tool, framework and job-title terms in English (PyTorch, MLOps, Machine Learning
Engineer) - translating them costs ATS keyword matches and reads as unidiomatic.]}

% ============================================================
%     CORE COMPETENCIES
% ============================================================

\section{ทักษะหลัก}
\vspace{1pt}
\begin{itemize}

\item \textbf{[SKILL_CATEGORY_1]}: [Specific skills, frameworks, tools.]

\item \textbf{[SKILL_CATEGORY_2]}: [Specific skills, frameworks, tools.]

\item \textbf{[SKILL_CATEGORY_3]}: [Specific skills, frameworks, tools.]

\item \textbf{[SKILL_CATEGORY_4]}: [Domain expertise, methods, approaches.]

\item \textbf{[SKILL_CATEGORY_5]}: [Tools, platforms, software.]

\end{itemize}

% ============================================================
%     PROFESSIONAL EXPERIENCE
% ============================================================

\section{ประสบการณ์ทำงาน}
\vspace{3pt}
\begin{itemize}

\needspace{5\baselineskip}
\item{\cventry{[YYYY-Present]}{[JOB_TITLE]}{[COMPANY]}{[CITY, COUNTRY]}{}{\vspace{1pt}
\begin{itemize}
    \item {[Achievement or responsibility 1 - be specific, use numbers where possible]}
    \item {[Achievement or responsibility 2]}
    \item {[Achievement or responsibility 3]}
    \item {[Achievement or responsibility 4]}
\end{itemize}}}

\vspace{3pt}

\needspace{5\baselineskip}
\item{\cventry{[YYYY-YYYY]}{[JOB_TITLE]}{[COMPANY]}{[CITY, COUNTRY]}{}{\vspace{1pt}
\begin{itemize}
    \item {[Achievement or responsibility 1]}
    \item {[Achievement or responsibility 2]}
    \item {[Achievement or responsibility 3]}
\end{itemize}}}

\vspace{3pt}

\needspace{5\baselineskip}
\item{\cventry{[YYYY-YYYY]}{[JOB_TITLE]}{[COMPANY]}{[CITY, COUNTRY]}{}{\vspace{1pt}
\begin{itemize}
    \item {[Achievement or responsibility 1]}
    \item {[Achievement or responsibility 2]}
\end{itemize}}}

\end{itemize}

% ============================================================
%     EDUCATION
% ============================================================

\section{การศึกษา}
\vspace{1pt}
\begin{itemize}

\needspace{5\baselineskip}
\item{\cventry{[YYYY-YYYY]}{[DEGREE] in [FIELD]}{[INSTITUTION]}{[CITY, COUNTRY]}{}{\vspace{1pt}
วิทยานิพนธ์: ``[THESIS_TITLE]'' [Brief description of research focus.]
}}

\vspace{3pt}

\needspace{5\baselineskip}
\item{\cventry{[YYYY-YYYY]}{[DEGREE] in [FIELD]}{[INSTITUTION]}{[CITY, COUNTRY]}{}{\vspace{1pt}
[Brief description or key topics.]
}}

\end{itemize}

% ============================================================
%     LANGUAGES
% ============================================================

\section{ภาษา}
\vspace{1pt}
\begin{itemize}
\item {ไทย ([LEVEL]), อังกฤษ ([LEVEL]), [LANGUAGE_3] ([LEVEL]).}
\end{itemize}

% ============================================================
%     PUBLICATIONS (optional)
% ============================================================

\section{ผลงานตีพิมพ์}
\vspace{1pt}
\begin{itemize}
\item {[AUTHORS] ([YEAR]). [TITLE]. [JOURNAL]. \href{[DOI_URL]}{DOI}}
\end{itemize}

% ============================================================
%     HONORS AND AWARDS (optional)
% ============================================================

\section{รางวัล}
\vspace{1pt}
\begin{itemize}
\item{\textbf{[AWARD_NAME]} -- [EVENT/ORGANIZATION] ([YEAR]).}
\end{itemize}

% ============================================================
%     REFERENCES
% ============================================================

\section{บุคคลอ้างอิง}
\vspace{1pt}
\begin{itemize}
\item{ยินดีให้ข้อมูลเมื่อได้รับการร้องขอ}
\end{itemize}

\end{document}
